\documentclass[10pt,a4paper]{article}

\usepackage[utf8]{inputenc}
\usepackage[T1]{fontenc}
\usepackage{lmodern}
\usepackage[a4paper,top=2.2cm,bottom=2cm,left=2.2cm,right=2.2cm]{geometry}
\usepackage[table]{xcolor}
\usepackage{tabularx}
\usepackage{booktabs}
\usepackage{array}
\usepackage{ragged2e}
\usepackage{fancyhdr}
\usepackage{enumitem}
\usepackage{microtype}
\usepackage{hyperref}
\usepackage{lastpage}
\usepackage{titlesec}

% Paleta sobria: un solo tono oscuro para texto/títulos, un gris para reglas y
% notas, y un acento discreto reservado únicamente para los tres niveles de
% riesgo que exige leerse a simple vista en la tabla.
\definecolor{ink}{HTML}{1F2933}
\definecolor{muted}{HTML}{6B7785}
\definecolor{rule}{HTML}{C7D0D8}
\definecolor{panel}{HTML}{F6F7F8}
\definecolor{accent}{HTML}{2E4756}
\definecolor{critical}{HTML}{8C2F2F}
\definecolor{high}{HTML}{A66A1E}
\definecolor{medium}{HTML}{7A7A2E}

\hypersetup{colorlinks=true,urlcolor=accent,linkcolor=accent,
  pdfauthor={Chiara Hernandez},pdftitle={Reporte ejecutivo de riesgos - Clínica privada}}
\setlength{\parindent}{0pt}
\setlength{\parskip}{4pt}
\setlist[itemize]{leftmargin=4.5mm,itemsep=1.5pt,topsep=2pt}
\renewcommand{\arraystretch}{1.3}

\pagestyle{fancy}
\fancyhf{}
\fancyhead[L]{\small\color{muted}Gestión de riesgos — Clínica privada}
\fancyhead[R]{\small\color{muted}Reporte ejecutivo}
\fancyfoot[L]{\small\color{muted}Chiara Hernandez \textbar\ LU 31964496 \textbar\ Comisión G}
\fancyfoot[R]{\small\color{muted}Página \thepage\ de \pageref{LastPage}}
\renewcommand{\headrulewidth}{0.4pt}
\renewcommand{\footrulewidth}{0pt}

\titleformat{\section}
  {\normalfont\Large\bfseries\color{ink}}
  {}{0em}{}
  [\vspace{1mm}{\color{rule}\titlerule[0.6pt]}]
\titlespacing*{\section}{0pt}{5mm}{3mm}

\newcommand{\decisionbox}[1]{%
\noindent\begin{tabularx}{\linewidth}{@{}p{0.2cm} X@{}}
\color{accent}\rule{0.6pt}{\baselineskip} & \small\textbf{Decisión requerida.}\ #1
\end{tabularx}\vspace{1mm}}

\newcommand{\riskcolor}[1]{%
  \ifnum\pdfstrcmp{#1}{Crítico}=0 critical\else
  \ifnum\pdfstrcmp{#1}{Alto}=0 high\else medium\fi\fi}

\newcommand{\planrow}[6]{%
\noindent\textbf{#1}\hfill \textit{0\% de avance}\par
\small #2\par
{\small\color{muted}Responsable: #3 \quad\textbar\quad Vencimiento: #4 \quad\textbar\quad Presupuesto: #5 \quad\textbar\quad Riesgo asociado: #6}
\par\vspace{2.5mm}}

\begin{document}

% --- Portada / encabezado ----------------------------------------------------
\thispagestyle{empty}
{\color{ink}\Huge\bfseries Reporte ejecutivo de riesgos}\par
\vspace{1mm}
{\color{muted}\large Clínica privada --- Registro inicial en SimpleRisk}\par
\vspace{2mm}
{\color{muted}\normalsize Dirigido al Directorio \textbar\ 5 de septiembre de 2026 \textbar\ Escenario ficticio}\par
\vspace{2mm}
{\color{rule}\rule{\linewidth}{0.8pt}}
\vspace{4mm}

\pagestyle{fancy}

\section*{Resumen ejecutivo}

La auditoría externa motivó la construcción de un registro inicial de riesgos para una clínica de 120 empleados que atiende aproximadamente 800 pacientes diarios y depende de historias clínicas digitales, datos de obras sociales y facturación.

El análisis identificó una exposición elevada: nueve de los diez riesgos evaluados requieren tratamiento prioritario o inmediato. La concentración de riesgos críticos en ransomware, accesos indebidos y phishing indica que la continuidad asistencial y la protección de datos sensibles dependen de fortalecer simultáneamente identidades, tecnología, procesos y capacidad de recuperación. Los controles actuales se consideran parciales y aún no existe evidencia de su eficacia sostenida. El riesgo restante, R09, corresponde a fallas físicas o ambientales y se mantiene en nivel medio bajo monitoreo.

\vspace{2mm}
\noindent
\begin{tabularx}{\linewidth}{XXXX}
\textbf{\large 10} & \textbf{\large 3} & \textbf{\large 6} & \textbf{\large 3} \\
{\scriptsize\color{muted}riesgos evaluados} & {\scriptsize\color{muted}críticos} & {\scriptsize\color{muted}altos} & {\scriptsize\color{muted}planes definidos}
\end{tabularx}

\vspace{3mm}
\decisionbox{Autorizar el inicio de los tres planes de acción ya definidos y exigir planes inmediatos adicionales para los riesgos críticos R01, R02 y R04. El costo inicial estimado de la cartera planificada es de USD 27.500.}

\section*{Top 5 riesgos por nivel}

\small
\begin{tabularx}{\linewidth}{>{\bfseries}c X >{\centering\arraybackslash}m{1.6cm} >{\centering\arraybackslash}m{1.5cm} X}
\toprule
\textbf{ID} & \textbf{Riesgo} & \textbf{P $\times$ I} & \textbf{Nivel} & \textbf{Consecuencia principal}\\
\midrule
R01 & Ransomware sobre historias clínicas y sistemas de atención & 4 $\times$ 5 = 20 & \textcolor{critical}{\textbf{Crítico}} & Paralización operativa, pérdida o exposición masiva de información.\\
R02 & Acceso indebido a historias clínicas & 4 $\times$ 4 = 16 & \textcolor{critical}{\textbf{Crítico}} & Vulneración de datos sensibles, consecuencias legales y pérdida de confianza.\\
R04 & Robo de credenciales mediante phishing & 4 $\times$ 4 = 16 & \textcolor{critical}{\textbf{Crítico}} & Acceso inicial para fraude, filtración o propagación de ataques.\\
R03 & Alteración de información clínica & 3 $\times$ 5 = 15 & \textcolor{high}{\textbf{Alto}} & Decisiones asistenciales basadas en datos incorrectos y pérdida de trazabilidad.\\
R07 & Imposibilidad de recuperar respaldos & 3 $\times$ 5 = 15 & \textcolor{high}{\textbf{Alto}} & Conversión de un incidente temporal en pérdida prolongada o permanente.\\
\bottomrule
\end{tabularx}
\normalsize

\vspace{2mm}
{\small\color{muted}Los empates se ordenaron priorizando la posible afectación de la atención y la continuidad de los servicios esenciales. La clasificación sigue la matriz de la cátedra: crítico 16--25, alto 10--15, medio 5--9, bajo 1--4.}

\newpage

% --- Página 2 -----------------------------------------------------------------
\section*{Estado de los planes de acción}

Los tres planes se encuentran en estado \textbf{Mitigation Planned}, con 0\% de avance. SimpleRisk mantiene por ahora el riesgo residual igual al inherente: la reducción solo podrá reconocerse una vez implementados y verificados los controles.

\planrow{PA01 --- Respaldos 3-2-1 y restauración}
{Copia aislada o inmutable, separación de credenciales y pruebas trimestrales. Primera prioridad porque reduce R07 y mejora la recuperación frente a R01.}
{Frank Castle --- Infraestructura y TI}{31/10/2026}{USD 8.000}{R07}

\planrow{PA03 --- Continuidad con obras sociales}
{Segundo enlace de conectividad, conmutación automática, reintentos y SLA. Protege admisión y facturación; requiere coordinación con proveedores.}
{Constanza Romero --- Administración y Finanzas}{15/11/2026}{USD 7.500}{R08}

\planrow{PA02 --- Integridad de historias clínicas}
{Trazabilidad, versionado y protección de cambios clínicos. Dirección Médica define los campos críticos; TI implementa los controles.}
{Juliana Gattas --- Dirección Médica}{15/12/2026}{USD 12.000}{R03}

\vspace{1mm}
{\small\color{muted}Inversión inicial estimada de la cartera: \textbf{\color{ink}USD 27.500}}

\section*{Hitos y criterios de aceptación}

\small
\begin{tabularx}{\linewidth}{>{\bfseries}m{1.1cm} X X}
\toprule
& \textbf{Primer hito} & \textbf{Criterio de cierre verificable}\\
\midrule
PA01 & Inventario de sistemas, RTO/RPO y arquitectura de respaldo aprobados. & Restauración integral trimestral dentro del tiempo objetivo, con evidencia y acciones correctivas.\\
PA03 & Segundo proveedor y cláusulas de disponibilidad seleccionados. & Prueba semestral de conmutación y recuperación de operaciones pendientes sin pérdida de registros.\\
PA02 & Campos clínicos críticos y eventos auditables definidos por Dirección Médica. & Registro protegido de usuario, fecha, valor anterior/nuevo, e historial consultable de modificaciones.\\
\bottomrule
\end{tabularx}
\normalsize

\section*{Gobernanza y seguimiento}

\begin{itemize}
  \item Presentar avance mensual al Comité de Dirección; escalar desvíos de fecha o presupuesto.
  \item El responsable de cada plan rinde cuentas por su ejecución; TI y Seguridad aportan evidencia técnica.
  \item Auditoría debe validar la eficacia de los controles antes de reducir el riesgo residual o declarar un plan completado.
  \item R01, R02 y R04 siguen sin un plan formal específico y requieren patrocinio inmediato del Directorio.
\end{itemize}

\decisionbox{Confirmar responsables, autorizar la ejecución de la cartera planificada y solicitar en un plazo de 30 días planes específicos para los tres riesgos críticos restantes.}

\newpage

% --- Página 3 -----------------------------------------------------------------
\section*{Recomendaciones prioritarias}

\begin{enumerate}[leftmargin=5.5mm,itemsep=3pt,label=\textbf{\arabic*.}]
  \item \textbf{Tratar inmediatamente los riesgos críticos.} Autorizar EDR, segmentación de red, MFA generalizado, revisión trimestral de privilegios, monitoreo de accesos y un programa continuo de concientización. El plan de respaldos mejora la recuperación, pero no reemplaza las medidas preventivas necesarias para R01, R02 y R04.

  \item \textbf{Ejecutar y controlar los tres planes ya definidos.} Mantener responsables, fechas y presupuesto; medir avances mensualmente y exigir evidencia de restauración, conmutación y trazabilidad antes de reducir el riesgo residual.

  \item \textbf{Implementar una estrategia de defensa en profundidad.} Superponer capas de control independientes: (i) identidad y autenticación multifactor; (ii) protección de endpoints, correo y red con monitoreo activo; y (iii) respaldos inmutables probados, continuidad operativa y un procedimiento de respuesta a incidentes. Si una capa falla, las restantes deben contener el daño.

  \item \textbf{Reforzar la protección frente a amenazas externas y suplantación.} El phishing que imita a una obra social o proveedor explota la confianza del personal para obtener credenciales o un punto de acceso inicial. Se recomienda MFA generalizado, validación por canal alternativo ante solicitudes sensibles, filtrado de correo y capacitación práctica basada en simulacros periódicos.

  \item \textbf{Proteger la infraestructura física y su ubicación.} La sala técnica debe estar separada de fuentes de riesgo ambiental (calor, agua, incendio) e incorporar sensores ambientales, detección de humo y filtraciones, UPS y generador probados periódicamente, y control de acceso físico.
\end{enumerate}

\section*{Hoja de ruta para decisión del Directorio}

\small
\begin{tabularx}{\linewidth}{>{\bfseries\color{ink}}m{2cm} X}
\toprule
0--30 días & Aprobar presupuesto y responsables; formalizar planes para R01, R02 y R04; activar MFA y revisión urgente de privilegios; definir RTO/RPO.\\
31--60 días & Implementar copia aislada/inmutable, segundo enlace de conectividad y primeras mejoras de auditoría clínica; ejecutar simulación de phishing y prueba de conmutación.\\
61--90 días & Probar restauración integral, verificar logs y versionado, revisar proveedores y reevaluar probabilidad, impacto y riesgo residual en SimpleRisk.\\
\bottomrule
\end{tabularx}
\normalsize

\vspace{3mm}
\decisionbox{El Directorio debe tratar la ciberseguridad como un riesgo asistencial y de negocio. La prioridad no es solo adquirir tecnología, sino demostrar que los controles funcionan mediante pruebas, métricas y revisión independiente.}

\vfill
{\scriptsize\color{muted}
\textbf{Fuentes de referencia:} Ley 25.326 de Protección de los Datos Personales; Ley 26.529 de Derechos del Paciente e Historia Clínica; CISA, \emph{StopRansomware Guide}; HHS, \emph{Health Industry Cybersecurity Practices}; NIST SP 800-53 Rev.\ 5 y NIST SP 800-34. Evaluación basada en la matriz de probabilidad por impacto proporcionada por la cátedra. Todos los datos personales y organizacionales utilizados son ficticios.}

\end{document}